\subsection{Why the results are what they are}
\label{sec:discussion-why}

The results have three separate origins, and the protocol measures each separately. The gap between the pages is a matter of counting. The default page counts a rating that defaulted whatever its position at the end of the year; the transition page counts only the ratings that are in a default category at the end, and a rating that defaults and is then withdrawn is in the withdrawals column. The speculative grades, where defaulted ratings are withdrawn at the end of the window in the case's lifetime check, therefore lose defaults from the transition page, and the loss is large exactly where the default rate is large. The data cannot separate this from informative withdrawal, because both remove would-be defaults from the numerator, and the real cohorts alone cannot say how much of the $\spLraDtwoCcc\%$ to $\spLraDfourCcc\%$ difference is which.

The bias of the treatments is a matter of selection. Proposition~\ref{prop:first-year} shows that a rating about to default is exposed, during the part of its last year that lies inside the cohort year, to a hazard of withdrawal $1+k$ times the baseline, so withdrawal is a competing event that hides defaults; Proposition~\ref{prop:later} shows that the same hazard acts a year ahead, so that the cohort of the next year is already short of the ratings about to default. The first effect is why removal and keeping are biased on the first cohort, and the second why the followed treatment is biased after it. Both are visible in the closed forms, and both are measured: the two-state closed form for the removal treatment at $k=9$ is within one standard error of the multi-state measurement.

The thin-grade results are a matter of discreteness. When the expected number of defaults in a grade is below one, the observed rate is zero in most samples, and an interval built on the observed rate cannot contain a positive PD (Proposition~\ref{prop:wald}); an interval that is centred elsewhere, as the Jeffreys and Agresti-Coull intervals are, contains the PD in the samples with no default and so reaches nominal coverage at the cost of a width of several hundred basis points around a PD of less than one.

\subsection{What a validator can do with them}
\label{sec:discussion-use}

Four practices follow from the measurements, each within what the data can support. A PD from agency statistics should carry the name of the page it was read from; the gap index of Equation~\eqref{eq:gap-index} and the risk-weight difference show what the alternative would have been. A treatment of withdrawn ratings should be chosen with the first cohort in view: where the end of every rating can be followed, as in the bank's own data, the followed treatment applied to a single cohort is unbiased in the model, and pooling cohorts reintroduces a selection bias that is a function of the informativeness. For a grade with an expected default count below about three, the Wald interval should not be reported; the Jeffreys interval covers, and its width states how little is known. And the matrix tests should not be taken as evidence against informative withdrawal: at levels of informativeness at which the PD is already wrong by a third they reject at rates of the order of their size.

The measurement does not say which page is right, or which value of $k$ holds for a given agency. It gives the range within which the PD moves, and the value of the range for capital is what a validator needs to size a margin of conservatism: at $k=9$ the removal treatment lowers the risk weight of S\&P's cohort mix by $\rwAbsremovedkNine$ points, which sets the order of magnitude for an adjustment of the risk weight if the withdrawals of the portfolio are suspected to be informative. A method that repairs the bias needs more than the transition page: it must anchor the informativeness on a second count of defaults, or report the PD as a function of $k$ with the value at which a decision flips. The engine of the case provides that family of estimates, and its evaluation is a separate study; the present measurement is what such an estimator has to beat.

\subsection{Negative and partial results}
\label{sec:discussion-negative}

Six results go against the simple story. The first is that the gap index of grade B at S\&P cannot be said to be below one half, so the strongest form of the claim holds for CCC to C only. The second is that following the withdrawn ratings, the treatment the guidelines ask for, is biased by $\followedBkNine\pm\followedBkNinese$ percent at $k=9$ once cohorts are pooled; its exactness on the first cohort does not carry over. The third is that the matrix tests do not see the withdrawals at $k=9$ and see them at only $\testmrkNineteenfive\%$ at $k=19$ with the better of the two. The fourth is that the first-order formula of the Theory section overstates the bias of grade CCC to C, where the PD is too large for a first-order expansion. The fifth is that the two-state closed form overstates the selection bias of the followed treatment on the multi-state generator ($\thBfollowedPoolkNine\%$ against $\followedBkNine\pm\followedBkNinese\%$), so the closed forms are a guide to the mechanism and not a substitute for the measurement on the grade at hand. The sixth is that the Jeffreys interval, the interval that behaves best at 50 obligors, is below nominal coverage in the thinnest grades at 1{,}000 obligors, because the enumeration of its coverage at a rare PD oscillates with $n$ and is not monotone.
